\documentclass[a4paper,10pt]{article}
\usepackage[top=25mm,bottom=24mm,left=27mm,right=27mm,headheight=15pt,headsep=8mm,footskip=12mm]{geometry}
\usepackage{fontspec,amsmath,amsthm,unicode-math}
\usepackage[spanish,es-nodecimaldot,es-noquoting]{babel}
\usepackage{microtype,xcolor,fancyhdr,titlesec,booktabs,longtable,array,tabularx}
\usepackage{graphicx,tikz,enumitem,float,needspace,fvextra}
\usetikzlibrary{arrows.meta,calc,positioning,decorations.pathreplacing,matrix}
\usepackage[unicode,hidelinks]{hyperref}
\usepackage[nameinlink,noabbrev]{cleveref}
\setmainfont{STIXTwoText-Regular.otf}[BoldFont=STIXTwoText-Bold.otf,ItalicFont=STIXTwoText-Italic.otf,BoldItalicFont=STIXTwoText-BoldItalic.otf]
\setmathfont{STIXTwoMath-Regular.otf}
\hypersetup{pdftitle={Estructura constitutiva del vacío electromagnético},pdfauthor={Oumar Haidara Fall; Rubén Ramos Balsa},pdfsubject={Artículo III. Borrador de trabajo para revisión académica}}
\newcommand{\Autores}{Oumar Haidara Fall \textperiodcentered\ Rubén Ramos Balsa}
\newcommand{\RR}{\mathbb R}\newcommand{\ZZ}{\mathbb Z}\newcommand{\QQ}{\mathbb Q}\newcommand{\CC}{\mathbb C}\newcommand{\FF}{\mathbb F}
\newcommand{\Id}{\operatorname{id}}\newcommand{\tr}{\operatorname{tr}}
\newcommand{\Span}{\operatorname{span}}\newcommand{\Cyl}{\operatorname{Cyl}}
\newcommand{\square}{\mdlgwhtsquare}
\newcommand{\TPK}{\mathrm{TPK}}\newcommand{\APP}{\mathrm{APP}}\newcommand{\TRIT}{\{-1,0,+1\}}
\newcommand{\HMT}{\mathrm{HMT}}\newcommand{\Tr}{\operatorname{Tr}}\newcommand{\dd}{\mathrm d}
\newtheorem{theorem}{Teorema}[section]
\newtheorem{lemma}[theorem]{Lema}
\newtheorem{proposition}[theorem]{Proposición}
\newtheorem{corollary}[theorem]{Corolario}
\theoremstyle{definition}\newtheorem{definition}[theorem]{Definición}\newtheorem{axiom}[theorem]{Axioma}
\theoremstyle{remark}\newtheorem{remark}[theorem]{Observación}
\numberwithin{equation}{section}
\setlength{\parindent}{1em}\setlength{\parskip}{3pt plus 1pt minus .5pt}
\setlength{\emergencystretch}{2em}
\renewcommand{\normalsize}{\fontsize{10.5}{13.4}\selectfont}\normalsize
\setlist{nosep,topsep=5pt}
\titleformat{\section}{\fontsize{13}{16}\selectfont\bfseries}{\thesection}{.65em}{}
\titleformat{\subsection}{\fontsize{11.2}{14}\selectfont\bfseries}{\thesubsection}{.65em}{}
\titleformat{\subsubsection}{\normalsize\bfseries}{\thesubsubsection}{.65em}{}
\AddToHook{cmd/section/before}{\clearpage}
\AddToHook{cmd/subsection/before}{\Needspace{7\baselineskip}}
\AddToHook{cmd/subsubsection/before}{\Needspace{6\baselineskip}}
\setcounter{tocdepth}{2}
\titlespacing*{\section}{0pt}{17pt}{9pt}\titlespacing*{\subsection}{0pt}{13pt}{6pt}
\pagestyle{fancy}\fancyhf{}
\fancyhead[L]{\fontsize{8}{10}\selectfont Estructura constitutiva del vacío electromagnético}
\fancyhead[R]{\fontsize{8}{10}\selectfont Artículo III}
\fancyfoot[C]{\thepage}\renewcommand{\headrulewidth}{.2pt}
\raggedbottom
\begin{document}
\begin{titlepage}
\centering\vspace*{28mm}
{\fontsize{23}{28}\selectfont\bfseries Estructura constitutiva\par del vacío electromagnético\par}
\vspace{12mm}
{\large\Autores\par}
\vspace{12mm}
{\normalsize Artículo III de la serie de Holografía Modular Triádica\par y Mecánica Dimensional\par}
\vfill
{\normalsize Borrador de trabajo para revisión académica\par}
\vspace{4mm}
{\small Integración y recopilación documental generadas por ChatGPT Astra.\par}
\vspace{18mm}
\end{titlepage}
\input{manuscrito/sections/00_resumen.tex}
\input{manuscrito/sections/01_introduccion.tex}
\clearpage\tableofcontents
\input{manuscrito/INCLUIR_NUCLEO_I_REV02.tex}
\input{manuscrito/INCLUIR_DEPENDENCIAS_II_REV02.tex}
\input{manuscrito/sections/03b_vacancias_prefactor.tex}
\input{manuscrito/sections/03_hojas_y_orientacion.tex}
\input{manuscrito/sections/04_respuesta_constitutiva.tex}
\input{manuscrito/sections/04b_reconstruccion_forma_lc.tex}
\input{manuscrito/sections/05_ciclo_catalan.tex}
\input{manuscrito/sections/06_pell_barbero.tex}
\input{manuscrito/sections/07_carga_y_cuantos.tex}
\input{manuscrito/sections/08_evaluacion_y_realizacion.tex}
\input{manuscrito/sections/09_conclusiones.tex}
\appendix
\input{manuscrito/sections/10_reproduccion.tex}
\input{manuscrito/sections/11_nucleo_finito.tex}
\input{manuscrito/bibliografia.tex}
\end{document}
